% Original analytic teaching diagram, 2026-11-12.
% Equal-weight mixture (1/2)N(mu1,1)+(1/2)N(mu2,1): (mu1,mu2) and (mu2,mu1) give the same law.
% Points (-1,2) and (2,-1) are mirror images across the diagonal mu1=mu2; shaded half-plane mu1<=mu2.
% Scale: 1 unit = 9 mm, window [-3,3]^2.
\definecolor{VB}{HTML}{7188A3}\definecolor{VBF}{HTML}{E5EAF0}
\definecolor{VG}{HTML}{829C8B}\definecolor{VGF}{HTML}{E8EEE9}
\definecolor{VY}{HTML}{B4A16F}\definecolor{VYF}{HTML}{F1ECDD}
\definecolor{VR}{HTML}{A7787D}\definecolor{VRF}{HTML}{F0E5E6}
\definecolor{VP}{HTML}{9588A5}\definecolor{VPF}{HTML}{ECE8F0}
\definecolor{VI}{HTML}{3B4252}\definecolor{VM}{HTML}{6B7280}\definecolor{VL}{HTML}{4C4D4F}
\begin{tikzpicture}[x=9mm,y=9mm,>=Stealth,font=\figurefont\fontsize{8}{10}\selectfont,
lab/.style={text=VM,font=\figurefont\fontsize{7.5}{9}\selectfont}]
\fill[VBF] (-3,-3) -- (3,3) -- (-3,3) -- cycle;
\draw[VL,->,line width=.5pt] (-3.2,0) -- (3.3,0) node[anchor=north east,text=VI] {$\mu_1$};
\draw[VL,->,line width=.5pt] (0,-3.2) -- (0,3.3) node[anchor=north east,text=VI] {$\mu_2$};
\foreach \t in {-2,-1,1,2} {\draw[VL,line width=.4pt] (\t,0) -- (\t,-0.1) node[below,lab] {$\t$};
  \draw[VL,line width=.4pt] (0,\t) -- (-0.1,\t) node[left,lab] {$\t$};}
\draw[VB,line width=1.1pt] (-3,-3) -- (3,3) node[anchor=north west,text=VB] {$\mu_1=\mu_2$};
\draw[VL,dashed,line width=.6pt] (-1,2) -- (2,-1);
\fill[VR] (-1,2) circle (0.11);
\fill[VR] (2,-1) circle (0.11);
\node[text=VR,anchor=south east] at (-1.05,2.05) {$(-1,2)$};
\node[text=VR,anchor=north west] at (2.05,-1.05) {$(2,-1)$};
\node[lab,anchor=west] at (-2.9,2.6) {$\mu_1\leq\mu_2$};
\end{tikzpicture}
